\documentclass[11pt]{article}
\usepackage[a4paper,margin=1in]{geometry}
\usepackage{fontspec}
\usepackage[boldfont=false]{xeCJK}
\setCJKmainfont{FandolSong-Regular.otf}
\usepackage{amsmath}
\usepackage{amssymb}
\usepackage{array}
\usepackage{booktabs}
\usepackage{graphicx}
\usepackage{adjustbox}
\usepackage{enumitem}
\setlist[itemize]{topsep=2pt,partopsep=0pt,parsep=0pt,itemsep=2pt,leftmargin=1.4em}
\setlist[enumerate]{topsep=2pt,partopsep=0pt,parsep=0pt,itemsep=2pt,leftmargin=1.6em}
\usepackage{parskip}
\usepackage{fvextra}
\fvset{breaklines=true,breakanywhere=true,fontsize=\small}
\setlength{\parindent}{0pt}

\begin{document}

\begin{center}
  {\LARGE\bfseries Financial information and decisions}\\[0.35em]
  {\large Igcse Business Studies}
\end{center}
\vspace{0.6em}

\section*{Why businesses need finance}

\textbf{Finance} 资金 means the money a business needs. A business needs finance to:

\begin{itemize}
\item start up — \textbf{start-up capital} 启动资金 to buy equipment and stock before it earns anything,

\item grow — money for \textbf{expansion} 扩张, such as new shops or machines,

\item run day to day — \textbf{working capital} 营运资金 to pay wages and suppliers while waiting for customers to pay,

\item survive a \textbf{cash-flow} 现金流 problem, when more money is going out than is coming in.

\end{itemize}

\subsection*{Sources of finance}

Finance can be \textbf{short-term} 短期 (paid back within a year) or \textbf{long-term} 长期 (over several years). It can come from inside the business (internal) or from outside (external).

\begin{center}
\adjustbox{max width=\linewidth}{%
\begin{tabular}{lll}
\toprule
\textbf{Source} & \textbf{Internal / external} & \textbf{Notes} \\
\midrule
\textbf{retained profit} 留存利润 & internal & profit kept in the business, not paid out \\
sale of \textbf{assets} 资产 & internal & selling things the business no longer needs \\
owners' savings & internal & the owners' own money \\
bank \textbf{loan} 贷款 & external & borrowed money, repaid with \textbf{interest} 利息 over time \\
bank \textbf{overdraft} 透支 & external & the bank lets the account go below zero, for short-term needs \\
issuing new \textbf{shares} 股份 & external & only for limited companies; raises money but shares out control \\
\textbf{debentures} 债券 & external & long-term loans to a company that pay interest \\
\textbf{leasing} 租赁 & external & renting equipment instead of buying it \\
\textbf{hire purchase} 分期付款 & external & buy now and pay in monthly parts \\
\textbf{trade credit} 贸易信贷 & external & pay suppliers later, for example after 30 days \\
\textbf{microfinance} 小额信贷 & external & small loans for tiny businesses in poorer areas \\
\textbf{crowdfunding} 众筹 & external & many people each give a small amount, often online \\
government \textbf{grants} 补助金 & external & money from the government that need not be repaid \\
\bottomrule
\end{tabular}}
\end{center}

\subsection*{Choosing a source}

The best source depends on how much money is needed and for how long, the cost (interest and fees), the size and legal type of the business (only companies can sell shares), and how much risk and control the owners will accept.

\section*{Cash flow and working capital}

\subsection*{Cash is not the same as profit}

\textbf{Cash} 现金 is the money a business can use right now. \textbf{Profit} 利润 is \textbf{revenue} 收入 minus costs over a period. A business can be making a profit yet still run out of cash — for example, if customers have not paid yet. Running out of cash can force even a profitable business to close.

\subsection*{The cash-flow forecast}

A \textbf{cash-flow forecast} 现金流预测 predicts the money coming in and going out each month. It uses:

\begin{itemize}
\item \textbf{inflows} 流入 — money coming in (sales, loans),

\item \textbf{outflows} 流出 — money going out (wages, rent, materials),

\item \textbf{net cash flow} 净现金流 — inflows minus outflows,

\item \textbf{opening balance} 期初余额 — the cash held at the start of the month,

\item \textbf{closing balance} 期末余额 — the cash held at the end of the month.

\end{itemize}

\[
\text{net cash flow} = \text{inflows} - \text{outflows}
\]

\[
\text{closing balance} = \text{opening balance} + \text{net cash flow}
\]

\subsection*{Working capital}

Working capital is the money available for day-to-day running. It is the \textbf{current assets} 流动资产 (cash and things soon turned into cash) minus the \textbf{current liabilities} 流动负债 (debts due within a year).

\[
\text{working capital} = \text{current assets} - \text{current liabilities}
\]

To overcome cash-flow problems, a business can arrange an overdraft, delay paying suppliers, ask customers to pay sooner, sell spare assets, or cut spending.

\section*{Income statement}

An \textbf{income statement} 损益表 shows how much profit a business made over a period, usually a year.

\begin{center}
\adjustbox{max width=\linewidth}{%
\begin{tabular}{ll}
\toprule
\textbf{Line} & \textbf{Meaning} \\
\midrule
revenue & money from sales \\
\textbf{cost of sales} 销售成本 & the cost of making or buying the goods that were sold \\
\textbf{gross profit} 毛利润 & revenue minus cost of sales \\
\textbf{expenses} 费用 & other running costs: rent, wages, advertising \\
\textbf{profit for the year} 年度利润 & gross profit minus expenses \\
\bottomrule
\end{tabular}}
\end{center}

\[
\text{gross profit} = \text{revenue} - \text{cost of sales}
\]

\[
\text{profit for the year} = \text{gross profit} - \text{expenses}
\]

Income statements are used by \textbf{stakeholders} 利益相关者: owners check the profit, lenders check whether their loans are safe, and managers look for problems.

\section*{Statement of financial position}

A \textbf{statement of financial position} 财务状况表 (also called a balance sheet) shows what a business owns and what it owes (its \textbf{liabilities} 负债) on one day.

\begin{itemize}
\item \textbf{non-current assets} 非流动资产 — things kept for the long term: buildings, machines.

\item \textbf{current assets} — short-term items: \textbf{inventory} 库存, money owed by customers, and cash.

\item \textbf{current liabilities} — debts due within a year: an overdraft, money owed to suppliers.

\item \textbf{non-current liabilities} 非流动负债 — debts due after more than a year: a long-term loan.

\item \textbf{owners' equity} 所有者权益 — the money the owners put in, plus profit kept in the business.

\end{itemize}

It is used to check what the business is worth, how much it owes, and whether it can pay its debts.

\section*{Analysis of accounts}

You can judge a business by working out \textbf{ratios} 比率 from its accounts. A ratio compares two figures.

\subsection*{Profitability ratios}

\textbf{Profitability} 盈利能力 ratios show how good a business is at making profit.

\[
\text{gross profit margin} = \frac{\text{gross profit}}{\text{revenue}} \times 100\%
\]

\[
\text{profit margin} = \frac{\text{profit for the year}}{\text{revenue}} \times 100\%
\]

\[
\text{ROCE} = \frac{\text{profit}}{\text{capital employed}} \times 100\%
\]

\begin{itemize}
\item \textbf{gross profit margin} 毛利率 — gross profit as a percentage of revenue.

\item \textbf{profit margin} 利润率 — profit for the year as a percentage of revenue.

\item \textbf{return on capital employed} 资本回报率 (ROCE) — profit as a percentage of the \textbf{capital employed} 所用资本 (money invested). It shows how well the investment is used.

\end{itemize}

A higher percentage is better for all three.

\subsection*{Liquidity ratios}

\textbf{Liquidity} 流动性 ratios show whether a business can pay its short-term debts.

\[
\text{current ratio} = \frac{\text{current assets}}{\text{current liabilities}}
\]

\[
\text{acid test ratio} = \frac{\text{current assets} - \text{inventory}}{\text{current liabilities}}
\]

\begin{itemize}
\item the \textbf{current ratio} 流动比率 compares current assets with current liabilities. Around 1.5 to 2 is healthy.

\item the \textbf{acid test ratio} 速动比率 is stricter: it leaves out inventory, because inventory is not yet sold.

\end{itemize}

\subsection*{Who uses account analysis}

\begin{itemize}
\item \textbf{internal users} (owners and managers) use it to control the business and to plan.

\item \textbf{external users} (banks, suppliers, the government and investors) use it to decide whether to lend, supply or invest.

\end{itemize}


\end{document}
